% GENERATED FILE. Edit canonical lessons, then run npm run generate:books.
% canonical-source-hash: fnv1a64:e573f860b52d90e5
% canonical-lessons: GU-C148-me-pidhi, GU-C148-tene, GU-C148-karyun, GU-C148-lidhun, GU-C148-did-practice

\chapter{What I did}
\label{ch:gu-what-i-did}
\hlchaptermodality{148}

\begin{chapteropening}
\textbf{By the end of this chapter:} \emph{I can say what I or someone else did, such as drank, ate, took, gave, wrote or said, with the past whose verb agrees with the object.}

Tell five past events side by side and say what each verb agrees with, then move one event to another doer.
\end{chapteropening}

\section[\texorpdfstring{\gu{મેં} … \gu{પીધી} (mẽ … pīdhī)}{મેં … પીધી (mẽ … pīdhī)}]{\gu{મેં} \gu{ચા} \gu{પીધી} — the past that agrees with the thing}

\label{lesson:GU-C148-me-pidhi}



\begin{quote}
Say \emph{I went} and \emph{she came}. Then three old words: \textbf{\gu{પીવું}}, \emph{to drink}; \textbf{\gu{ચા}}, \emph{tea}; \textbf{\gu{પાણી}}, \emph{water}.
\end{quote}

\begin{grammarlens}[title={Grammar Lens: \gu{હું} becomes \gu{મેં}}]
\textbf{\gu{હું} \gu{ગયો}} keeps \emph{I} in its plain form. A verb that does something \emph{to} something, like drinking tea or eating rice, changes that in the past: the doer is marked, and \textbf{\gu{હું}} becomes \textbf{\gu{મેં}} (\emph{mẽ}).

\begin{quote}
\textbf{\gu{મેં} \gu{ચા} \gu{પીધી}.} — \emph{mẽ cā pīdhī.} — I drank tea.
\end{quote}

\textbf{\gu{મેં}} is the same for a man and a woman. The verb has stopped listening to the drinker.
\end{grammarlens}

\subsection*{You'll want to know: the verb follows the tea}
\textbf{\gu{પીધી}} ends in \textbf{-\gu{ી}} because \textbf{\gu{ચા}} is feminine. Change the drink and the verb changes with it:

\begin{itemize}
\raggedright
  \item \textbf{\gu{મેં} \gu{ચા} \gu{પીધી}.} — \emph{mẽ cā pīdhī} — tea, feminine;
  \item \textbf{\gu{મેં} \gu{પાણી} \gu{પીધું}.} — \emph{mẽ pāṇī pīdhũ} — water, neuter;
  \item \textbf{\gu{મેં} \gu{દૂધ} \gu{પીધું}.} — \emph{mẽ dūdh pīdhũ} — milk, neuter.
\end{itemize}

This is the agreement \textbf{\gu{મારી} \gu{ચા}} showed you: the ending follows the thing. English has nothing like it, so it is worth saying slowly a few times.

\subsection*{Your turn}
\begin{itemize}
\raggedright
  \item \emph{Say it:} \emph{mẽ cā pīdhī}
  \item \emph{Say it:} \emph{mẽ dūdh pīdhũ}
  \item \emph{Say it:} \textbf{\gu{મેં} \gu{તમારી} \gu{ચા} \gu{પીધી}!} — \emph{mẽ tamārī cā pīdhī!}, I drank your tea!
  \item \emph{Contrast:} \emph{hũ gayo} and \emph{mẽ cā pīdhī}: which verb follows the doer, and which the drink
\end{itemize}

\begin{tcolorbox}[breakable,colback=teal!4,colframe=teal!35!black,title={Before you move on}]
What does \textbf{\gu{હું}} become before a past verb with an object? (\textbf{\gu{મેં}}.) What does that verb agree with? (The object: \textbf{\gu{ચા}} gives \textbf{\gu{પીધી}}.)
\end{tcolorbox}
\section[\texorpdfstring{\gu{મેં} · \gu{તેં} · \gu{તેણે} · \gu{અમે} … (mẽ · tẽ · teṇe …)}{મેં · તેં · તેણે · અમે … (mẽ · tẽ · teṇe …)}]{\gu{તેણે}, \gu{તેમણે} — who did it}

\label{lesson:GU-C148-tene}



\begin{quote}
Say \emph{I drank tea} and \emph{I drank water}, and say why the two verbs end differently.
\end{quote}

\begin{grammarlens}[title={Grammar Lens: one person}]
\noindent
\begin{tabularx}{\linewidth}{@{}>{\raggedright\arraybackslash}X>{\raggedright\arraybackslash}X>{\raggedright\arraybackslash}X@{}}
\toprule
\textbf{who} & \textbf{plain} & \textbf{in the object past} \\
\midrule
I & \gu{હું} & \gu{મેં} \\
you (familiar) & \gu{તું} & \gu{તેં} \\
he, she & \gu{તે} & \gu{તેણે} \\
\bottomrule
\end{tabularx}

\begin{quote}
\textbf{\gu{તેણે} \gu{ચા} \gu{પીધી}.} — \emph{teṇe cā pīdhī.} — He drank tea. She drank tea.
\end{quote}

\textbf{\gu{તેણે}} is the form CDIAL records for Gujarati under the old Sanskrit stem \emph{ta-} (CDIAL 5612), the stem behind \textbf{\gu{તે}} and \textbf{\gu{તેઓ}}.
\end{grammarlens}

\begin{grammarlens}[title={Grammar Lens: more than one}]
\noindent
\begin{tabularx}{\linewidth}{@{}>{\raggedright\arraybackslash}X>{\raggedright\arraybackslash}X>{\raggedright\arraybackslash}X@{}}
\toprule
\textbf{who} & \textbf{plain} & \textbf{in the object past} \\
\midrule
we & \gu{અમે} & \gu{અમે} \\
you (respectful) & \gu{તમે} & \gu{તમે} \\
they & \gu{તેઓ} & \gu{તેમણે} \\
\bottomrule
\end{tabularx}

\textbf{\gu{અમે}} and \textbf{\gu{તમે}} do not change at all: \textbf{\gu{અમે} \gu{ચા} \gu{પીધી}}, \textbf{\gu{તમે} \gu{પાણી} \gu{પીધું}}. Only \textbf{\gu{તેઓ}} moves, to \textbf{\gu{તેમણે}}, built on the \textbf{\gu{તેમ}-} of \textbf{\gu{તેમને}}.
\end{grammarlens}

\subsection*{Your turn}
Say these aloud:

\begin{itemize}
\raggedright
  \item \emph{teṇe pāṇī pīdhũ}
  \item \textbf{\gu{તમે} \gu{ચા} \gu{પીધી}?} — \emph{tame cā pīdhī?}, did you drink tea?
  \item \textbf{\gu{તેમણે} \gu{દૂધ} \gu{પીધું}} — \emph{temṇe dūdh pīdhũ}
\end{itemize}

\begin{tcolorbox}[breakable,colback=teal!4,colframe=teal!35!black,title={Before you move on}]
Which two doers keep their plain form? (\textbf{\gu{અમે}} and \textbf{\gu{તમે}}.) What does \textbf{\gu{તેઓ}} become? (\textbf{\gu{તેમણે}}.)
\end{tcolorbox}
\section[\texorpdfstring{\gu{કર્યું} · \gu{આપ્યું} · \gu{જોયું} … (karyũ · āpyũ …)}{કર્યું · આપ્યું · જોયું … (karyũ · āpyũ …)}]{\gu{કર્યું} — the regular past: stem and -\gu{્યું}}

\label{lesson:GU-C148-karyun}



\begin{quote}
Say \emph{she drank tea} and \emph{they drank milk}. Then \textbf{\gu{કામ} \gu{કરવું}}, \emph{to work}, \emph{to do work}.
\end{quote}

\begin{grammarlens}[title={Grammar Lens: -\gu{્યું} on the stem}]
Most verbs build this past in one move: take the stem and add \textbf{-\gu{્યું}}, a virama and then \textbf{\gu{યું}}, so the stem's last consonant joins the \textbf{\gu{ય}}.

\begin{itemize}
\raggedright
  \item \textbf{\gu{કર}-} $\to$ \textbf{\gu{કર્યું}} — \emph{karyũ} — did (the \textbf{\gu{ર}} rides on top as a small hook);
  \item \textbf{\gu{આપ}-} $\to$ \textbf{\gu{આપ્યું}} — \emph{āpyũ} — gave;
  \item \textbf{\gu{લખ}-} $\to$ \textbf{\gu{લખ્યું}} — \emph{lakhyũ} — wrote;
  \item \textbf{\gu{કહ}-} $\to$ \textbf{\gu{કહ્યું}} — \emph{kahyũ} — said;
  \item \textbf{\gu{જો}-} $\to$ \textbf{\gu{જોયું}} — \emph{joyũ} — saw (a stem ending in a vowel takes \textbf{\gu{યું}} with no virama).
\end{itemize}

\begin{quote}
\textbf{\gu{મેં} \gu{કામ} \gu{કર્યું}.} — \emph{mẽ kām karyũ.} — I did the work. \textbf{\gu{તેણે} \gu{પત્ર} \gu{લખ્યો}.} — \emph{teṇe patra lakhyo.} — She wrote a letter. \textbf{\gu{મેં} \gu{તેને} \gu{પૈસા} \gu{આપ્યા}.} — \emph{mẽ tene paisā āpyā.} — I gave him the money. \textbf{\gu{તેણે} \gu{મને} \gu{પત્ર} \gu{આપ્યો}.} — \emph{teṇe mane patra āpyo.} — She gave me a letter. \textbf{\gu{મેં} \gu{બે} \gu{ફોન} \gu{કર્યા}.} — \emph{mẽ be phon karyā.} — I made two phone calls.
\end{quote}

The four endings again: \textbf{-\gu{્યો}}, \textbf{-\gu{ી}}, \textbf{-\gu{્યું}}, \textbf{-\gu{્યા}}. \textbf{\gu{પત્ર}} is masculine, so \textbf{\gu{લખ્યો}}; \textbf{\gu{પૈસા}} is plural, so \textbf{\gu{આપ્યા}}. In the feminine the \textbf{\gu{ય}} drops: \textbf{\gu{કરી}}, \textbf{\gu{આપી}}, \textbf{\gu{જોઈ}}.
\end{grammarlens}

\subsection*{Your turn}
Say these aloud:

\begin{itemize}
\raggedright
  \item \emph{mẽ kām karyũ}
  \item \emph{teṇe patra lakhyo}
  \item \emph{mẽ tene paisā āpyā}
  \item \textbf{\gu{તેમણે} \gu{શું} \gu{કહ્યું}?} — \emph{temṇe śũ kahyũ?}, what did they say?
\end{itemize}

\begin{tcolorbox}[breakable,colback=teal!4,colframe=teal!35!black,title={Before you move on}]
What is added to a stem like \textbf{\gu{કર}-} to make this past? (\textbf{-\gu{્યું}}.) Why \textbf{\gu{લખ્યો}} in \emph{she wrote a letter}? (\textbf{\gu{પત્ર}} is masculine.)
\end{tcolorbox}
\section[\texorpdfstring{\gu{લીધું} · \gu{ખાધું} · \gu{પીધું} (līdhũ · khādhũ · pīdhũ)}{લીધું · ખાધું · પીધું (līdhũ · khādhũ · pīdhũ)}]{\gu{લીધું}, \gu{ખાધું}, \gu{પીધું} — three pasts in -\gu{ધું}}

\label{lesson:GU-C148-lidhun}



\begin{quote}
Say \emph{I did the work} and \emph{she wrote a letter}. Then \textbf{\gu{લેવું}}, \emph{to take}, and \textbf{\gu{ખાવું}}, \emph{to eat}.
\end{quote}

\begin{grammarlens}[title={Grammar Lens: -\gu{ધું}}]
\textbf{\gu{પીધું}} has been doing this since the first lesson of the chapter. Three short, common verbs share it: instead of \textbf{-\gu{્યું}} they take \textbf{-\gu{ધું}}.

\begin{itemize}
\raggedright
  \item \textbf{\gu{લેવું}} $\to$ \textbf{\gu{લીધું}} — \emph{līdhũ} — took;
  \item \textbf{\gu{ખાવું}} $\to$ \textbf{\gu{ખાધું}} — \emph{khādhũ} — ate;
  \item \textbf{\gu{પીવું}} $\to$ \textbf{\gu{પીધું}} — \emph{pīdhũ} — drank.
\end{itemize}

Their endings are the usual four: \textbf{\gu{ખાધો}}, \textbf{\gu{ખાધી}}, \textbf{\gu{ખાધું}}, \textbf{\gu{ખાધા}}.

\begin{quote}
\textbf{\gu{મેં} \gu{ભાત} \gu{ખાધો}.} — \emph{mẽ bhāt khādho.} — I ate rice. \textbf{\gu{તેણે} \gu{રોટલી} \gu{ખાધી}.} — \emph{teṇe roṭlī khādhī.} — She ate a roṭlī. \textbf{\gu{મેં} \gu{ટિકિટ} \gu{લીધી}.} — \emph{mẽ ṭikiṭ līdhī.} — I took a ticket. \textbf{\gu{તેણે} \gu{ભાત} \gu{લીધો}.} — \emph{teṇe bhāt līdho.} — He took some rice. \textbf{\gu{તેણે} \gu{દવા} \gu{લીધી}.} — \emph{teṇe davā līdhī.} — He took the medicine.
\end{quote}
\end{grammarlens}

\subsection*{Your turn}
Say these aloud:

\begin{itemize}
\raggedright
  \item \emph{mẽ bhāt khādho}
  \item \emph{teṇe roṭlī khādhī}
  \item \emph{mẽ ṭikiṭ līdhī}
\end{itemize}

\begin{tcolorbox}[breakable,colback=teal!4,colframe=teal!35!black,title={Before you move on}]
Which three verbs take \textbf{-\gu{ધું}}? (\textbf{\gu{લેવું}}, \textbf{\gu{ખાવું}}, \textbf{\gu{પીવું}}.) Why \textbf{\gu{ખાધો}} with \textbf{\gu{ભાત}}? (\textbf{\gu{ભાત}} is masculine.)
\end{tcolorbox}
\section[Practice]{What I did — a practice}

\label{lesson:GU-C148-did-practice}



\begin{quote}
Say \emph{I ate rice}, then \emph{I did the work}.

\begin{itemize}
\raggedright
  \item \emph{Write it:} \textbf{\gu{ડ}} once from memory
\end{itemize}
\end{quote}

\begin{tcolorbox}[breakable,skin=enhanced,colback=orange!6,colframe=orange!45!black,boxrule=0.5pt,arc=1mm,left=6pt,right=6pt,top=4pt,bottom=4pt,fonttitle=\bfseries,title={Read this}]
\begin{quote}
\textbf{\gu{ગઈકાલે} \gu{હું} \gu{બજાર} \gu{ગઈ}.} \textbf{\gu{મેં} \gu{કેરી} \gu{લીધી}.} \textbf{\gu{સાંજે} \gu{મારો} \gu{ભાઈ} \gu{આવ્યો}.} \textbf{\gu{તેણે} \gu{ભાત} \gu{ખાધો} \gu{અને} \gu{ચા} \gu{પીધી}.} \textbf{\gu{અમે} \gu{તેને} \gu{પુસ્તક} \gu{આપ્યું}.}
\end{quote}

Say these aloud:

\begin{itemize}
\raggedright
  \item the five lines aloud
  \item for each verb, who or what its ending agrees with
\end{itemize}

\textbf{\gu{ગઈ}} and \textbf{\gu{આવ્યો}} follow whoever moved. \textbf{\gu{લીધી}}, \textbf{\gu{ખાધો}}, \textbf{\gu{પીધી}} and \textbf{\gu{આપ્યું}} follow the thing: the mango, the rice, the tea, the book.
\end{tcolorbox}

\subsection*{Your turn}
\begin{itemize}
\raggedright
  \item \emph{Answer:} \textbf{\gu{તમે} \gu{ગઈકાલે} \gu{શું} \gu{ખાધું}?} — \emph{tame gaīkāle śũ khādhũ?}
\end{itemize}

\begin{tcolorbox}[breakable,colback=teal!4,colframe=teal!35!black,title={Before you move on}]
In \textbf{\gu{મેં} \gu{કેરી} \gu{લીધી}}, why \textbf{-\gu{ી}}? (\textbf{\gu{કેરી}} is feminine.) In \textbf{\gu{હું} \gu{બજાર} \gu{ગઈ}}, why \textbf{-\gu{ઈ}}? (The speaker is a woman.)
\end{tcolorbox}
